%
% vox-music-player.tex
%
% Z Specification for the voxd Music Player's view state machine.
%
% The player owns NO playback machinery. Its play/stop/advance transitions are
% StartRadio / RadioOff / RadioRotate on a single-album selection, already
% fuzz-clean and ProB-checked in docs/audio-programs.tex. This model pins the one
% new, stable thing: the PlayerView --- the projection of voxd's single active
% source onto a lux scene --- and its invariants across the player transitions:
%
%   I1  at most one album playing
%   I2  now-playing present iff playing
%   I3  a played album is catalogued (and a playing album cannot be removed)
%
% The phase-2 connection / subscribe / lease lifecycle depends on lux's
% not-yet-pinned subscribe API and is deferred to an amendment (see section 6).
%

\documentclass[a4paper,10pt,fleqn]{article}
\usepackage[margin=1in]{geometry}
\usepackage{fuzz}
\usepackage[colorlinks=true,linkcolor=blue,citecolor=blue,urlcolor=blue]{hyperref}

\begin{document}

\title{The voxd Music Player: A Z Specification}
\author{Formal model of the PlayerView state machine}
\date{July 2026}
\hypersetup{
  pdftitle={The voxd Music Player: A Z Specification},
  pdfauthor={Punt Labs},
  pdfsubject={Z Specification},
  pdfcreator={fuzz/probcli}
}
\maketitle

\section{Introduction}

The Music Player is a headless app inside \texttt{voxd} that projects the
daemon's one active source onto a lux scene and translates in-scene clicks back
into playback commands. It adds no playback machinery: playing a saved album is
\texttt{ProgramService.replay\_album}, stopping is \texttt{off}, and advancing
within the album is the playback loop driving \texttt{SelectionPlayback.rotate}.
Those transitions --- \texttt{StartRadio}, \texttt{RadioOff},
\texttt{RadioRotate}, and the \texttt{MusicRemove} refusal of a playing album ---
are already modelled and \texttt{fuzz}-clean in \texttt{docs/audio-programs.tex}.

This specification pins the one new, stable object: the \emph{PlayerView}, and
the three invariants the design must preserve across every player transition.
The phase-2 lux connection/subscribe/lease lifecycle depends on lux's
not-yet-pinned subscribe API and is deferred (Section~\ref{sec:deferred}).

\section{Basic Types and Modes}

An \texttt{ALBUM} is an opaque catalog identifier --- the id \texttt{vox music
list} prints and the scene's Play button carries. It is a catalog key, never a
filesystem path; the model treats it as opaque, exactly as
\texttt{audio-programs.tex} does.

\begin{zed}
  [ALBUM]
\end{zed}

The player is a two-mode machine: nothing playing, or one album playing.

\begin{zed}
  PlayerMode ::= idle | playing
\end{zed}

\section{State}
\label{sec:state}

The \texttt{PlayerView} carries the catalog of saved albums, the mode, the
optional playing album (a set of at most one, mirroring the ``optional Part''
cursor of \texttt{audio-programs.tex}), and the now-playing cursor
\texttt{(index, of)} --- the 1-based track position and the album's track count,
both zero when idle.

The predicate carries every invariant. \texttt{I1} bounds the playing album to
at most one. \texttt{I2} is the biconditional the design names ``now-playing
present iff playing'': the mode is \texttt{playing} exactly when one album is
playing, exactly when the cursor is live (\texttt{of} $\geq 1$). \texttt{I3}
contains the playing album in the catalog. The cursor-validity conjuncts tie the
index into $1 \upto of$ while playing and pin it to zero while idle.

\begin{schema}{PlayerView}
  catalogued : \power ALBUM \\
  mode : PlayerMode \\
  album : \power ALBUM \\
  index : \nat \\
  of : \nat
  \where
  \# album \leq 1 \\
  album \subseteq catalogued \\
  (mode = playing \iff \# album = 1) \\
  (mode = idle \iff album = \emptyset) \\
  (mode = playing \implies (1 \leq index \land index \leq of \land of \geq 1)) \\
  (mode = idle \implies (index = 0 \land of = 0))
\end{schema}

The now-playing cursor is the \texttt{(index, of)} pair the status API reports as
``Part $N$ of $M$''; \texttt{I2} makes ``present iff playing'' a fact of the state
schema, not a runtime check. Because \texttt{PlayerMode} has only two
inhabitants, $mode = idle$ and $mode = playing$ are complementary, so the two
biconditionals are consistent: $album = \emptyset \iff \# album = 0 \iff mode
\neq playing$.

\section{Initialisation}

The daemon starts with a (possibly empty) catalog and nothing playing.

\begin{schema}{PlayerInit}
  PlayerView' \\
  cat0? : \power ALBUM
  \where
  catalogued' = cat0? \\
  mode' = idle \\
  album' = \emptyset \\
  index' = 0 \\
  of' = 0
\end{schema}

\section{Player Transitions}

Each transition is the projection of an \texttt{audio-programs.tex} operation
onto the view. They are gated on mode alone, never on a session --- any client
(the lux button, the CLI, or an MCP tool) may drive any of them, which is why the
scene must re-push on every state change regardless of origin.

\subsection{PlayerPlay --- play a saved album}

A client plays the album \texttt{target?}, whose ready-track count is
\texttt{count?}. The album must be catalogued and non-empty ($count? \geq 1$):
\texttt{replay\_album} refuses an unknown id and an album with no playable
tracks. There is \emph{no} mode precondition: from \texttt{idle} this is
\texttt{StartRadio}; while already \texttt{playing} it is
\texttt{SwitchSelection}, displacing the old album. Either way the after-state
has exactly one album playing, so \texttt{I1} holds and two albums never play at
once.

\begin{schema}{PlayerPlay}
  \Delta PlayerView \\
  target? : ALBUM \\
  count? : \nat
  \where
  target? \in catalogued \\
  count? \geq 1 \\
  catalogued' = catalogued \\
  mode' = playing \\
  album' = \{ target? \} \\
  index' = 1 \\
  of' = count?
\end{schema}

The cursor starts at track~1 of \texttt{count?}. Because $count? \geq 1$, the
after-state satisfies the cursor-validity invariant ($1 \leq 1 \leq of'$ and
$of' \geq 1$), and because $target? \in catalogued$, \texttt{I3}
($album' \subseteq catalogued'$) holds. \texttt{PlayerPlay} is the single verb
for both ``start'' and ``switch'', exactly as \texttt{StartRadio} $\defs$
\texttt{SwitchSelection} in \texttt{audio-programs.tex}.

\subsection{PlayerTrackEnd --- auto-advance within the album}

When the current track ends, the playback loop advances the cursor to another
track of the same album, avoiding an immediate repeat when the album holds more
than one and replaying when it holds exactly one. This is \texttt{RadioRotate}
projected onto the cursor: the album and its count are unchanged, only the index
moves.

\begin{schema}{PlayerTrackEnd}
  \Delta PlayerView
  \where
  mode = playing \\
  catalogued' = catalogued \\
  mode' = playing \\
  album' = album \\
  of' = of \\
  1 \leq index' \land index' \leq of \\
  (of \geq 2 \implies index' \neq index) \\
  (of = 1 \implies index' = index)
\end{schema}

The automatic track-end advance and a user-driven \texttt{next}/\texttt{skip} are
one transition, mirroring \texttt{RadioNext} $\defs$ \texttt{RadioRotate}. The
mode, album, and count are all held, so \texttt{I1}, \texttt{I2}, and \texttt{I3}
are trivially preserved; only the cursor moves, and it stays in range.

\subsection{PlayerStop --- stop, return to idle}

Stopping tears the source down and returns the view to \texttt{idle}: no album,
no cursor. This is \texttt{RadioOff}. It is enabled while playing; a stop while
already idle is a no-op the daemon need not model as a transition.

\begin{schema}{PlayerStop}
  \Delta PlayerView
  \where
  mode = playing \\
  catalogued' = catalogued \\
  mode' = idle \\
  album' = \emptyset \\
  index' = 0 \\
  of' = 0
\end{schema}

The after-state is the \texttt{idle} shape: \texttt{I2} holds because emptying
the album empties the cursor, discharging ``stop returns to idle''.

\section{Catalog Transitions}

The catalog changes under the player via \texttt{music new} and \texttt{music
remove}. They are modelled here only for their interaction with the view ---
specifically, that a playing album cannot be removed.

\subsection{CatalogAdd --- a fresh album appears}

\texttt{music new} files a fresh album. It grows the catalog and leaves playback
untouched, so the scene re-pushes with a new row but the same now-playing.

\begin{schema}{CatalogAdd}
  \Delta PlayerView \\
  new? : ALBUM
  \where
  new? \notin catalogued \\
  catalogued' = catalogued \cup \{ new? \} \\
  mode' = mode \\
  album' = album \\
  index' = index \\
  of' = of
\end{schema}

Because the catalog only grows, $album \subseteq catalogued'$ still holds and
\texttt{I3} is preserved with playback unchanged.

\subsection{CatalogRemove --- delete a non-playing album}

\texttt{music remove} deletes a catalogued album, but the playing album is
refused --- the model image of the \texttt{active\_backing\_locators} guard and of
\texttt{MusicRemove}'s precondition in \texttt{audio-programs.tex}. The
precondition $gone? \notin album$ is exactly ``you cannot delete the album you
are playing''.

\begin{schema}{CatalogRemove}
  \Delta PlayerView \\
  gone? : ALBUM
  \where
  gone? \in catalogued \\
  gone? \notin album \\
  catalogued' = catalogued \setminus \{ gone? \} \\
  mode' = mode \\
  album' = album \\
  index' = index \\
  of' = of
\end{schema}

The guard is what preserves \texttt{I3} across a removal: with $gone? \notin
album$ and $album \subseteq catalogued$, removing \texttt{gone?} leaves
$album \subseteq catalogued'$. Deleting the playing album would break
$album' \subseteq catalogued'$, so the invariant is not decoration --- it forces
the guard rather than the guard being an ad-hoc check.

\section{Deferred: the connection / subscribe / lease lifecycle}
\label{sec:deferred}

Phase~2 adds the receive leg: a persistent WebSocket subscription over the public
\texttt{LuxRestClient} extension, decoding \texttt{music.play} / \texttt{music.stop}
into the transitions above, and a ``Music'' menu callback held by a $\sim$30s
lease renewed by contact. That lifecycle --- \texttt{disconnected} $\to$
\texttt{connected} $\to$ \texttt{subscribed} $\to$ \texttt{lease-renewed}, with a
click received while disconnected dropped and the lease expiring unless renewed
--- is itself a stateful machine with invariants (at most one live subscription;
no dispatch while disconnected). It is \emph{not} modelled here because its
transitions depend on lux's subscribe/lease API, which is not yet pinned. When
lux pins that API, this specification gains a \texttt{Connection} schema and its
operations as an amendment, before the phase-2 subscribe loop is implemented.

The separation is deliberate and mirrors the design: the \texttt{PlayerView} is
stable today because a transition's \emph{trigger} --- a lux click or a CLI call
--- is immaterial to the view's invariants. Only the plumbing that carries the
trigger is unpinned.

\end{document}
